% ==============================================================================
% 12_part3_results.tex — ส่วนที่ 3: ผลการประเมินตนเองและแผนพัฒนาหลักสูตร
% ==============================================================================

\section*{ผลการประเมินตนเองรายเกณฑ์ (Summary of Self-Assessment Scores)}
\addcontentsline{toc}{section}{ผลการประเมินตนเองรายเกณฑ์}

\subsection*{3.1 ผลการประเมินตนเองรายข้อ (Detailed Sub-criteria Assessment)}

ตารางต่อไปนี้แสดงผลการประเมินตนเองรายข้อย่อยตามเกณฑ์ AUN-QA Version~4.0 ทั้ง 8 เกณฑ์ โดยเครื่องหมาย $\checkmark$ แสดงระดับคะแนนที่หลักสูตรประเมินตนเองได้ในแต่ละข้อ (คะแนน~1--7)

\begingroup
\scriptsize
\setlength{\tabcolsep}{2pt}
\renewcommand{\arraystretch}{1.2}
\begin{longtable}{|p{0.6cm}|>{\raggedright\arraybackslash}p{9.2cm}|c|c|c|c|c|c|c|}
  \hline
  \rowcolor{gray!20}
  & \textbf{Criteria / Requirements} & \textbf{1} & \textbf{2} & \textbf{3} & \textbf{4} & \textbf{5} & \textbf{6} & \textbf{7} \\
  \hline
  \endfirsthead
  \hline
  \rowcolor{gray!20}
  & \textbf{Criteria / Requirements (ต่อ)} & \textbf{1} & \textbf{2} & \textbf{3} & \textbf{4} & \textbf{5} & \textbf{6} & \textbf{7} \\
  \hline
  \endhead
  \hline
  \multicolumn{9}{r}{\footnotesize (ต่อหน้าถัดไป)} \\
  \endfoot
  \hline
  \endlastfoot

  % ── Criterion 1 ──────────────────────────────────────────────────────────────
  \multicolumn{9}{|l|}{\cellcolor{gray!15}\textbf{1\quad Expected Learning Outcomes}} \\
  \hline
  1.1 & The programme to show that the expected learning outcomes are appropriately formulated in accordance with an established learning taxonomy, are aligned to the vision and mission of the university, and are known to all stakeholders. & & $\checkmark$ & & & & & \\
  \hline
  1.2 & The programme to show that the expected learning outcomes for all courses are appropriately formulated and are aligned to the expected learning outcomes of the programme. & & $\checkmark$ & & & & & \\
  \hline
  1.3 & The programme to show that the expected learning outcomes consist of both generic outcomes (related to written and oral communication, problem-solving, information technology, teambuilding skills, etc) and subject specific outcomes (related to knowledge and skills of the study discipline). & & $\checkmark$ & & & & & \\
  \hline
  1.4 & The programme to show that the requirements of the stakeholders, especially the external stakeholders, are gathered, and that these are reflected in the expected learning outcomes. & & $\checkmark$ & & & & & \\
  \hline
  1.5 & The programme to show that the expected learning outcomes are achieved by the students by the time they graduate. & & $\checkmark$ & & & & & \\
  \hline

  % ── Criterion 2 ──────────────────────────────────────────────────────────────
  \multicolumn{9}{|l|}{\cellcolor{gray!15}\textbf{2\quad Programme Structure and Content}} \\
  \hline
  2.1 & The specifications of the programme and all its courses are shown to be comprehensive, up-to-date, and made available and communicated to all stakeholders. & & $\checkmark$ & & & & & \\
  \hline
  2.2 & The design of the curriculum is shown to be constructively aligned with achieving the expected learning outcomes. & & $\checkmark$ & & & & & \\
  \hline
  2.3 & The design of the curriculum is shown to include feedback from stakeholders, especially external stakeholders. & & $\checkmark$ & & & & & \\
  \hline
  2.4 & The contribution made by each course in achieving the expected learning outcomes is shown to be clear. & & $\checkmark$ & & & & & \\
  \hline
  2.5 & The curriculum to show that all its courses are logically structured, properly sequenced, and are integrated. & & $\checkmark$ & & & & & \\
  \hline
  2.6 & The curriculum to have option(s) for students to pursue major and/or minor specialisations. & & $\checkmark$ & & & & & \\
  \hline
  2.7 & The programme to show that its curriculum is reviewed periodically following an established procedure and that it remains up-to-date and relevant to industry. & & $\checkmark$ & & & & & \\
  \hline

  % ── Criterion 3 ──────────────────────────────────────────────────────────────
  \multicolumn{9}{|l|}{\cellcolor{gray!15}\textbf{3\quad Teaching and Learning Approach}} \\
  \hline
  3.1 & The teaching and learning activities are shown to be constructive, relevant, and aligned with the intended learning outcomes, and reflect the educational philosophy of the institution. & & $\checkmark$ & & & & & \\
  \hline
  3.2 & The teaching and learning activities are shown to provide opportunities for students to participate actively and take responsibility for their own learning. & & $\checkmark$ & & & & & \\
  \hline
  3.3 & The teaching and learning activities are shown to promote students actively engaged in the learning process. & & $\checkmark$ & & & & & \\
  \hline
  3.4 & The teaching and learning activities are shown to support and encourage students to be innovative and have creative thought and entrepreneurial mindset. & & $\checkmark$ & & & & & \\
  \hline
  3.5 & The teaching and learning activities are shown to promote lifelong learning. & & $\checkmark$ & & & & & \\
  \hline
  3.6 & The teaching and learning activities are shown to be regularly reviewed and evaluated for their effectiveness and improvement in aligning with the expected learning outcomes. & & $\checkmark$ & & & & & \\
  \hline

  % ── Criterion 4 ──────────────────────────────────────────────────────────────
  \multicolumn{9}{|l|}{\cellcolor{gray!15}\textbf{4\quad Student Assessment}} \\
  \hline
  4.1 & The student assessment is shown to use a variety of assessment methods that are constructively aligned with the intended learning outcomes. & & $\checkmark$ & & & & & \\
  \hline
  4.2 & The programme to show that its student assessment policies and regulations are explicit, fair, and consistently applied. & & $\checkmark$ & & & & & \\
  \hline
  4.3 & The programme to show that the student progression and degree completion standards are explicit, communicated, and consistently applied. & & $\checkmark$ & & & & & \\
  \hline
  4.4 & The programme to show that its student assessment uses rubrics, marking schemes, and regulations that ensure validity, reliability, and fairness. & & $\checkmark$ & & & & & \\
  \hline
  4.5 & The programme to show that student assessment is used to measure both the achievement of the intended course learning outcomes and the contribution of each course to the programme learning outcomes. & & $\checkmark$ & & & & & \\
  \hline
  4.6 & The programme to show that the results of the student assessment are used in a timely manner for providing feedback to help students improve their learning. & & $\checkmark$ & & & & & \\
  \hline
  4.7 & The programme to show that its student assessment methods are subject to continuous review in order to ensure their validity, reliability, and relevance to the student assessment policy. & & $\checkmark$ & & & & & \\
  \hline

  % ── Criterion 5 ──────────────────────────────────────────────────────────────
  \multicolumn{9}{|l|}{\cellcolor{gray!15}\textbf{5\quad Academic Staff}} \\
  \hline
  5.1 & The programme to show that there is an academic staff planning process that aligns to the needs of the programme. & & $\checkmark$ & & & & & \\
  \hline
  5.2 & The programme to show that the workload of academic staff is measured, monitored and managed to ensure quality in teaching, research and service. & & $\checkmark$ & & & & & \\
  \hline
  5.3 & The programme to show that the competences of academic staff are defined, assessed, and communicated. & & $\checkmark$ & & & & & \\
  \hline
  5.4 & The programme to show that the workload and responsibilities of academic staff are allocated in accordance with the qualifications, experience, and skill sets. & & $\checkmark$ & & & & & \\
  \hline
  5.5 & The programme to show that the criteria and procedures for promotion of academic staff are clearly defined, communicated, and based on merit. & & $\checkmark$ & & & & & \\
  \hline
  5.6 & The programme to show that the rights, privileges, benefits, relationships, and responsibilities of academic staff are clearly defined and communicated. & & $\checkmark$ & & & & & \\
  \hline
  5.7 & The programme to show that there is a staff training and development plan that is developed based on training needs analysis. & & $\checkmark$ & & & & & \\
  \hline
  5.8 & The programme to show that there is a systematic approach in managing and evaluating academic staff performance that includes reward and recognition. & & $\checkmark$ & & & & & \\
  \hline

  % ── Criterion 6 ──────────────────────────────────────────────────────────────
  \multicolumn{9}{|l|}{\cellcolor{gray!15}\textbf{6\quad Student Support Services}} \\
  \hline
  6.1 & The programme to show that the student intake policy, criteria, and procedures are clearly defined, communicated, publicly available, and up-to-date. & & $\checkmark$ & & & & & \\
  \hline
  6.2 & The programme to show that there are adequate short-term and long-term plans for student support services. & & $\checkmark$ & & & & & \\
  \hline
  6.3 & The programme to show that there is a system to monitor, respond and report on student progress, academic performance, and workload. & & $\checkmark$ & & & & & \\
  \hline
  6.4 & The programme to show that there are co-curricular activities, student competitions, and other support services to enhance the learning experience and employability of students. & & $\checkmark$ & & & & & \\
  \hline
  6.5 & The programme to show that the competences of support staff are defined, assessed and communicated. & & $\checkmark$ & & & & & \\
  \hline
  6.6 & The programme to show that student support services are evaluated and benchmarked for improvement. & & $\checkmark$ & & & & & \\
  \hline

  % ── Criterion 7 ──────────────────────────────────────────────────────────────
  \multicolumn{9}{|l|}{\cellcolor{gray!15}\textbf{7\quad Facilities and Infrastructure}} \\
  \hline
  7.1 & The programme to show that the physical resources are adequate to support the learning and research of students. & & $\checkmark$ & & & & & \\
  \hline
  7.2 & The programme to show that laboratory equipment and facilities are adequate, up-to-date, readily available, and effectively deployed. & & $\checkmark$ & & & & & \\
  \hline
  7.3 & The programme to show that there is a digital library with adequate and up-to-date resources and facilities that are in line with ICT advancement. & & $\checkmark$ & & & & & \\
  \hline
  7.4 & The programme to show that the IT systems adequately support teaching, learning, research and service of academic staff, support staff and students. & & $\checkmark$ & & & & & \\
  \hline
  7.5 & The programme to show that the computer and network infrastructure is adequate and accessible to academic staff and students throughout the institution. & & $\checkmark$ & & & & & \\
  \hline
  7.6 & The programme to show that it complies with environment, health, safety and accessibility standards in the institution. & & $\checkmark$ & & & & & \\
  \hline
  7.7 & The programme to show that the physical, social, and psychological environment is conducive to education and research as well as student well-being. & & $\checkmark$ & & & & & \\
  \hline
  7.8 & The programme to show that the competences of the facilities and infrastructure support staff are defined, assessed and communicated. & & $\checkmark$ & & & & & \\
  \hline
  7.9 & The programme to show that the quality of facilities and infrastructure is evaluated and benchmarked for improvement. & & $\checkmark$ & & & & & \\
  \hline

  % ── Criterion 8 ──────────────────────────────────────────────────────────────
  \multicolumn{9}{|l|}{\cellcolor{gray!15}\textbf{8\quad Output and Outcomes}} \\
  \hline
  8.1 & The programme to show that the pass rate, dropout rate, and average time to graduate are established, monitored, and benchmarked for improvement. & & $\checkmark$ & & & & & \\
  \hline
  8.2 & The programme to show that the employability, self-employment, entrepreneurship and further study of graduates are established, monitored, and benchmarked for improvement. & & $\checkmark$ & & & & & \\
  \hline
  8.3 & The programme to show that research and creative work outputs of students and academic staff are established, monitored, and benchmarked for improvement. & & $\checkmark$ & & & & & \\
  \hline
  8.4 & The programme to show that the achievement of programme outcomes by students is established and monitored for improvement. & & $\checkmark$ & & & & & \\
  \hline
  8.5 & The programme to show that the satisfaction of stakeholders is established, monitored, and benchmarked for improvement. & & $\checkmark$ & & & & & \\
  \hline
\end{longtable}
\endgroup

\begin{flushleft}
  \scriptsize \textbf{หมายเหตุ}: คะแนน~1 = ไม่มีหลักฐาน, คะแนน~2 = มีหลักฐานบางส่วน, คะแนน~3 = มีหลักฐานชัดเจน, คะแนน~4 = มีระบบที่ดี, คะแนน~5 = มีระบบที่ดีมากและวัดได้, คะแนน~6 = ดีเลิศ, คะแนน~7 = เป็นแบบอย่าง ($\checkmark$ = ระดับคะแนนที่ประเมินตนเอง)
\end{flushleft}

\subsection*{3.2 สรุปผลการประเมินตนเองรายเกณฑ์}

ตาราง~\ref{tbl:sar_summary} แสดงผลการประเมินตนเองรายเกณฑ์ของหลักสูตรวิศวกรรมศาสตรมหาบัณฑิต สาขาวิชาวิศวกรรมคอมพิวเตอร์และปัญญาประดิษฐ์ มหาวิทยาลัยราชภัฏอุตรดิตถ์ ปีการศึกษา~2568

\begin{table}[H]
  \centering
  \caption{สรุปผลการประเมินตนเอง AUN-QA Version~4.0 ปีการศึกษา~2568}
  \label{tbl:sar_summary}
  \begingroup
  \scriptsize
  \setlength{\tabcolsep}{3pt}
  \renewcommand{\arraystretch}{1.15}
  \begin{tabularx}{\linewidth}{|c|>{\raggedright\arraybackslash}X|c|>{\raggedright\arraybackslash}X|}
    \hline
    \textbf{Cri.} & \textbf{เกณฑ์} & \textbf{คะแนน} & \textbf{เหตุผลหลัก} \\
    \hline
    1 & Expected Learning Outcomes & 2 & PLO ครบ Bloom's Taxonomy (C3--C6), SH Survey n=31 + External SHs~9 แห่ง, PLO Achievement Plan ครบ 5 PLO \\
    \hline
    2 & Programme Structure and Content & 2 & มคอ.2 + มคอ.3 ครบถ้วน, OBE/BCD + Constructive Alignment, CLO-PLO Mapping ครบ 33 วิชา \\
    \hline
    3 & Teaching and Learning Approach & 2 & IEL ใน มคอ.3 ทุกวิชา, Active Learning หลากหลาย, Guest Lectures + Community Engagement, Cross-faculty Input \\
    \hline
    4 & Student Assessment & 2 & Rubric 4 ระดับทุก CLO, PLO Achievement Criteria ครบ 5 PLO, ระบบอุทธรณ์, Constructive Alignment ครบ \\
    \hline
    5 & Academic Staff & 2 & คณาจารย์~10 คน ปร.ด./วศ.ด. ทุกคน (สมอ.08), URU-PSF, Workload System, ทุนวิจัย~3 โครงการ + IPITEx Bronze, Training Records~3 คน \\
    \hline
    6 & Student Support Service & 2 & SmartURU + Advisor System, แผนระยะสั้น/ยาว~5 แผน, ผลประเมินบริการ~4.51/5.0 \\
    \hline
    7 & Facilities and Infrastructure & 2 & Computer Lab~277 เครื่อง, WiFi~795 จุด, ฐานข้อมูล~8 ฐาน, ผลประเมิน~4.34/5.0 \\
    \hline
    8 & Output and Outcomes & 2 & Pass Rate~100\%, Research~7 รายการ (Scopus Q2~×2, TCI กลุ่ม~1, ทุน~3, IPITEx Bronze) = 1.00 ชิ้น/คน/ปี, PLO~100\% ครบ, Satisfaction~4.94/5.00 \\
    \hline
    \multicolumn{2}{|l|}{\textbf{คะแนนเฉลี่ยรวม}} & \textbf{2.00} & \\
    \hline
  \end{tabularx}
  \endgroup
\end{table}

\vspace{0.5em}
\textbf{หมายเหตุ}: คะแนน~1 = ไม่บรรลุ (หลักสูตรใหม่ ปีแรก), คะแนน~2 = บรรลุบางส่วน (มีระบบและหลักฐาน แต่ยังพัฒนาได้), คะแนน~3 = บรรลุ (มีระบบครบและมีผลลัพธ์ดี), คะแนน~4--7 สำหรับหลักสูตรที่มีข้อมูลย้อนหลัง 3--5 ปี

\subsection*{จุดแข็งของหลักสูตร (Programme Strengths)}

\begin{enumerate}
  \item \textbf{คณาจารย์คุณภาพสูง}: อาจารย์ประจำหลักสูตร~10 คน (ตาม สมอ.08 มีผล ภาค~2/2568) ทุกคนมีตำแหน่งทางวิชาการและวุฒิปริญญาเอก ความเชี่ยวชาญตรงกับ PLOs และมีผลงานวิจัยที่ได้รับทุนสนับสนุนจากแหล่งภายนอก

  \item \textbf{การเชื่อมโยงกับชุมชนและภาคอุตสาหกรรม}: มี External Stakeholders~9 หน่วยงานที่มีส่วนร่วมจริง ทั้งการวิจัยร่วม การอบรม การส่งบุคลากรเรียน และการรับโจทย์วิจัยจริงมาใช้ในวิทยานิพนธ์

  \item \textbf{โครงสร้างพื้นฐานดิจิทัลที่แข็งแกร่ง}: WiFi~795 จุด ห้องปฏิบัติการคอมพิวเตอร์~277 เครื่อง ฐานข้อมูลวิจัยนานาชาติ~7 ฐาน และระบบ IT สนับสนุนการศึกษา~10 ระบบ

  \item \textbf{การออกแบบหลักสูตรตาม OBE/IEL}: ทุก มคอ.3 ระบุ IEL ใน CLO ทุกข้อ Constructive Alignment ชัดเจน และมีสัปดาห์ Community Workshop กำหนดไว้ในตาราง

  \item \textbf{ผลงานวิจัยโดดเด่นในปีแรก}: อาจารย์มีผลงาน~7 รายการ (Scopus Q2 จำนวน~2 บทความ, TCI กลุ่ม~1, ทุน Fundamental Fund/สวก./CBR, รางวัล IPITEx Bronze Award) นักศึกษาได้รับรางวัลระดับภาคเหนือและสามารถเป็นวิทยากรสอน AI ให้ครูโรงเรียน~3 แห่ง
\end{enumerate}

\subsection*{จุดที่ต้องพัฒนา (Areas for Improvement)}

\begin{enumerate}
  \item \textbf{การเก็บข้อมูล Evidence อย่างเป็นระบบ}: Training Records รายบุคคลอาจารย์, Formal Monitoring Records นักศึกษา~4 คน --- รายงานการประชุมยกร่างหลักสูตร (\hyperlink{ref:AUNQA-1-2}{AUNQA-1-2}) จัดทำเป็นรายงาน LaTeX แล้ว (report\_meeting\_2566-09-15)

  \item \textbf{CLO/PLO Achievement Data จริง}: ต้องรวบรวม CLO Achievement จาก มคอ.5 ภาค~2/2568 ทุกวิชา และพัฒนา PLO Aggregation Framework เพื่อรายงาน PLO Achievement Level ระดับหลักสูตร

  \item \textbf{Benchmarking รอบต่อไป}: C6 กำหนด Formal Benchmarking Partners ภายนอก (มรภ.เชียงราย/เชียงใหม่/ลำปาง) แล้ว รอดำเนินการเปรียบเทียบจริงในภาค~1/2569; C8 มี Baseline Benchmark แล้ว (R8.1/8.3/8.5) ต้องพัฒนา Employer Satisfaction Survey เมื่อบัณฑิตรุ่นแรกสำเร็จการศึกษา

  \item \textbf{Training Needs Analysis (C5)}: จัดทำ TNA อย่างเป็นทางการจากการสำรวจคณาจารย์รายบุคคล และพัฒนา Individual Development Plans

  \item \textbf{Safety Standards for Research Lab (C7)}: พัฒนา ESPReL หรือมาตรฐานความปลอดภัยห้องปฏิบัติการวิจัยเฉพาะสำหรับ CPE\&AI

\end{enumerate}

\subsection*{แผนพัฒนาหลักสูตร (Action Plan 2569)}

\begin{table}[H]
  \centering
  \caption{แผนพัฒนาหลักสูตร ปีการศึกษา~2569}
  \label{tbl:action_plan}
  \begingroup
  \footnotesize
  \setlength{\tabcolsep}{3pt}
  \renewcommand{\arraystretch}{1.2}
  \begin{tabularx}{\linewidth}{|c|>{\raggedright\arraybackslash}X|>{\raggedright\arraybackslash}X|l|>{\raggedright\arraybackslash}X|r|>{\raggedright\arraybackslash}X|}
    \hline
    \textbf{Cri.} & \textbf{โครงการ} & \textbf{กลุ่มเป้าหมาย} & \textbf{ระยะเวลา} & \textbf{ผู้รับผิดชอบ} & \textbf{งบ (บาท)} & \textbf{KPI} \\
    \hline
    4 & รวบรวม CLO Achievement + จัดทำ PLO Aggregation Report & นักศึกษาปี~2568 ($n=4$) & สิ้นปีการศึกษา 2568 & อาจารย์ผู้สอนทุกวิชา & --- (ภาระงาน) & PLO Aggregation Report ครอบคลุม PLO 5 ตัว ส่งภายใน 31 พ.ค. 2569 \\
    \hline
    6 & Benchmarking ขยายผล (จาก Baseline ปัจจุบัน: มช., สจล., มรภ.นครปฐม) ไป KPI เชิงสัดส่วน Retention/Satisfaction/Time-to-graduate & 3 หลักสูตร ป.โท CE/AI ที่ baseline แล้ว & ภาค 1/2569 & คณะกรรมการบริหารหลักสูตร & 15{,}000 & รายงาน Benchmark KPI 1 ฉบับ ภายใน 31 ต.ค. 2569 \\
    \hline
    8 & Formal Employer Survey ระดับ ป.โท (เตรียมการล่วงหน้าก่อนบัณฑิตรุ่นแรกสำเร็จการศึกษา 2569--2570) + Stakeholder Satisfaction รอบ 2 & สถานประกอบการ/นายจ้างศักยภาพ $\geq$5 แห่ง + นักศึกษา 4 คน + External SHs & ภาค 1/2569 & คณะกรรมการบริหารหลักสูตร & 12{,}000 & แบบสำรวจ Employer ตอบกลับ $\geq$5 ราย + รายงาน Satisfaction $\geq$4.5/5.00 \\
    \hline
  \end{tabularx}
  \endgroup
\end{table}
